\section{Introduction}
\label{sec:introduction}

Artificial intelligence has historically been engineered as a largely static artifact. A model is trained, a policy is selected, a prompt is written, an architecture is deployed, and the resulting system is then evaluated as if its intelligence were a property fixed at release time. This view is convenient for benchmarking and for product deployment, yet it is increasingly misaligned with the systems now emerging around large language models, tool-using agents, automated program synthesis, and evolutionary search. Contemporary AI systems do not merely map inputs to outputs. They can inspect execution traces, write and execute code, critique their own reasoning, revise their prompts, call external tools, store episodic memories, generate new training tasks, and propose modifications to the very software that implements them. These capabilities suggest a shift from static AI to \emph{self-evolving AI}: systems that improve through feedback-driven loops over their own components.

This survey studies that shift. We use the term \emph{AI Self-Evolution} to denote a family of methods in which an AI system modifies some part of itself---its prompts, memory, tools, agent code, architecture, data curriculum, or neural weights---as a result of feedback collected from interaction with tasks, users, environments, verifiers, or other agents. The field is not a single algorithm. It is a convergence zone where language-agent reasoning, AutoML and neural architecture search, evolutionary computation, reinforcement learning from verifiable feedback, program synthesis, and open-ended science automation are beginning to share concepts. A prompt-refinement loop such as Self-Refine \cite{madaan2023selfrefine}, a memory-based verbal reinforcement learner such as Reflexion \cite{shinn2023reflexion}, a symbolic-backpropagation agent optimizer \cite{zhou2024agent}, a meta-agent that writes new agent architectures \cite{hu2024automated}, a Darwinian archive of self-modifying coding agents \cite{zhang2025darwin}, a runtime monkey-patching agent inspired by the G{\"o}del machine \cite{yin2024godel,schmidhuber2003godel}, an evolutionary coding system such as AlphaEvolve \cite{novikov2025alphaevolve}, and a zero-data self-play reasoner \cite{zhao2025absolute} all differ in mechanism. Yet they can be compared under the same abstraction: a system observes performance, converts feedback into a modification proposal, validates the proposal, and conditionally retains the modified system.

The timing of this convergence is not accidental. Classical AI already contained the ingredients: Turing's behavioral framing of machine intelligence \cite{turing1950computing}, recursive self-improvement thought experiments, G{\"o}del-machine theory \cite{schmidhuber2003godel}, meta-learning, evolutionary algorithms, and AutoML. What changed in the 2020s was the availability of foundation models that can operate over the representations in which modern AI systems are built. Prompts are text, tool descriptions are text, code is text, benchmark reports are text, stack traces are text, scientific papers are text, and even many architectural decisions can be serialized as text. Large language models therefore act as general-purpose variation operators: they can propose a new prompt, edit a Python class, design a test, summarize a failure, synthesize a training example, or explain why an earlier attempt failed. When such proposals are coupled to automated evaluators---unit tests, code executors, human preference models, software-engineering benchmarks, mathematical verifiers, web task success criteria, or peer-review rubrics---the result is a practical loop of self-modification and selection.

The research challenge is to understand when these loops constitute genuine self-evolution rather than ordinary adaptation or offline optimization. A system that is merely retrained by engineers after deployment is not self-evolving in the sense used here. A model that passively receives more data under a fixed update rule may be online learning, but not necessarily self-evolution. Conversely, a language agent that keeps its weights fixed but revises its memory and prompts after every failure can be self-evolving at the representational level at which its behavior is actually controlled. The key issue is not whether the modified component is neural, symbolic, or procedural. The key issue is whether the AI system participates in the generation, evaluation, and retention of modifications to its own operational machinery.

This chapter introduces the background, definition, boundaries, and contributions of the survey. Section~\ref{subsec:intro-background} traces the path from static AI and the Turing test to G{\"o}del machines, AutoML, evolutionary architecture search, large language models, and modern self-improving agents. Section~\ref{subsec:intro-definition} provides a formal definition of AI Self-Evolution as a feedback-driven modification process over system components. Section~\ref{subsec:intro-related} distinguishes self-evolution from continual learning, online learning, self-supervised learning, meta-learning, AutoML, reinforcement learning, and conventional agent memory. Section~\ref{subsec:intro-contributions} summarizes the survey's contributions: a Five Evolution Loops framework, analysis of more than nineteen open-source and research systems, a cross-domain unification of AutoML/NAS/LLM/evolutionary-computation/agent research, and a researcher-network map of the field. Section~\ref{subsec:intro-structure} outlines the rest of the paper.

\subsection{Background: From Static AI to Self-Evolving AI}
\label{subsec:intro-background}

\subsubsection{The static-system assumption in early artificial intelligence}

The earliest public imagination of artificial intelligence was already entangled with questions of evaluation and change. Turing's 1950 proposal replaced the metaphysical question ``Can machines think?'' with an operational imitation game in which a machine's intelligence is judged through interaction \cite{turing1950computing}. The test did not prescribe a specific learning mechanism, but it established a durable pattern: intelligence would be inferred from behavior under a task protocol. Most subsequent benchmarks inherited this framing. A system is trained or programmed, placed behind an interface, and judged by its outputs. The protocol may be conversational, mathematical, perceptual, strategic, or embodied, but the object under test is normally treated as fixed during evaluation.

This static-system assumption shaped the organization of modern machine learning. A training phase produces parameters $\theta$; a validation phase chooses hyperparameters; a test phase estimates generalization. The deployed model is then represented as a function $f_{\theta}:\mathcal{X}\rightarrow\mathcal{Y}$, and the central research question is how to learn $\theta$ from a dataset $D$. Even reinforcement learning, which emphasizes interaction, usually freezes the learned policy when reporting benchmark scores. This division is scientifically useful because it separates search from assessment, reduces leakage, and makes comparisons reproducible. Yet it also narrows the idea of an AI system to the final artifact rather than the process that produced and continues to alter that artifact.

Symbolic AI and expert systems were also mostly static after design time. Engineers encoded rules, taxonomies, production systems, or search strategies. A knowledge base could be updated, but the update was typically performed by humans or by a narrow acquisition routine. The expert system did not redesign its inference engine, rewrite its knowledge-acquisition workflow, or choose a new representational scheme in response to failures. The system's intelligence was in its manually specified structure. Machine learning later replaced hand-written rules with learned functions, but the deployed artifact remained a snapshot.

The static assumption was reinforced by software-engineering practice. A model release is versioned, audited, and monitored. If performance degrades, developers collect data, retrain the model, patch code, or adjust prompts. The system does not usually hold authority to alter those components itself. This separation between the object being improved and the human process that improves it is the main boundary that AI Self-Evolution begins to erode. In a self-evolving system, some part of the improvement process becomes internal to the system's own operation.

The historical importance of static evaluation should not be understated. Without frozen benchmarks, it would be impossible to measure progress in speech recognition, computer vision, machine translation, game playing, or code generation. However, the same practice can obscure a new class of systems whose competence resides not only in their current policy but also in their ability to generate better future policies. A static snapshot of an evolving agent is analogous to a photograph of a learning organism: useful, but incomplete. For self-evolving AI, we must evaluate both the current behavior and the evolution operator that changes that behavior over time.

\subsubsection{Self-reference and the G{\"o}del-machine ideal}

Long before current language agents, researchers considered whether an intelligent machine could improve its own learning algorithm. Schmidhuber's G{\"o}del machine is the canonical formal proposal \cite{schmidhuber2003godel}. It describes a self-referential system that contains a representation of its own software and a proof searcher. The machine rewrites itself only after proving that the rewrite will improve expected utility according to a formal objective. The theoretical beauty of the G{\"o}del machine lies in its treatment of self-modification as a rational action. It does not merely mutate itself; it requires a proof that the mutation is beneficial.

The original G{\"o}del-machine ideal is far stronger than most modern systems. Current agents rarely possess complete formal models of their environments, their utility functions, or their own implementation. They cannot generally prove that a code patch will improve long-run utility. Yet the idea introduced several principles that reappear in contemporary self-evolution. First, the system must have some access to a description of itself. Second, self-modification must be evaluated relative to an objective. Third, the system must decide whether to accept or reject candidate modifications. Fourth, the improvement process is recursive: the mechanism that modifies the system may itself be modified.

Modern systems relax formal proof into empirical validation. Instead of proving that a modification is globally utility-improving, they run tests, benchmarks, simulations, or external evaluations. G{\"o}del Agent, for example, explicitly invokes the self-referential inspiration of the G{\"o}del machine but implements runtime self-modification through Python monkey patching and validation on tasks \cite{yin2024godel}. Darwin G{\"o}del Machine combines the self-referential ambition with Darwinian search: a foundation model proposes code-level changes to agents, and an archive retains variants that improve on benchmarks \cite{zhang2025darwin}. These systems do not deliver the mathematical guarantees of Schmidhuber's construction. They instead trade proof for empirical selection, scalability, and compatibility with messy software environments.

This shift from proof to evaluation is one of the defining compromises of practical self-evolution. Formal proof is attractive but brittle in open domains. Empirical evaluation is noisy but scalable. The self-evolving systems surveyed in this paper almost always use the latter. Reflexion stores verbal feedback after failed episodes \cite{shinn2023reflexion}; Self-Refine iteratively critiques and revises outputs \cite{madaan2023selfrefine}; Agent Symbolic Learning produces language gradients over prompts, tools, and pipelines \cite{zhou2024agent}; ADAS evaluates newly generated agent programs \cite{hu2024automated}; AlphaEvolve uses automated evaluators to select program variants \cite{novikov2025alphaevolve}; Absolute Zero uses a code executor as a verifier for self-generated tasks \cite{zhao2025absolute}. The unifying pattern is not proof-carrying self-rewrite but feedback-carrying self-modification.

G{\"o}del-machine theory also clarifies why self-evolution is potentially unstable. If a system can modify its own modification procedure, it may improve faster than a system limited to a fixed update rule. But the same recursion can amplify errors, reward hacking, unsafe objectives, or distributional drift. A self-modification that improves a short benchmark may degrade robustness; a prompt revision that increases success on one task may remove safety constraints; a code patch that passes tests may exploit the evaluator. Thus self-evolution inherits both the promise and the danger of self-reference. The survey's emphasis on loops, validation, rollback, and governance follows directly from this tension.

\subsubsection{From hand-designed learning to AutoML and neural architecture search}

A second lineage comes from automating the design of learning systems. AutoML asks whether the pipeline that turns data into a model can itself be optimized: data preprocessing, feature engineering, architecture selection, hyperparameter tuning, training schedules, and ensembling. Neural architecture search (NAS) narrows the question to architectures, often using reinforcement learning, evolutionary algorithms, differentiable relaxations, or weight sharing. Methods such as NASNet \cite{zoph2018learning}, Efficient NAS \cite{pham2018efficient}, DARTS \cite{liu2018darts}, and regularized evolution for image classifiers \cite{real2019regularized} showed that search can discover architectures competitive with human designs.

AutoML and NAS matter for self-evolution because they shift intelligence from a learned model to a search process over models. The artifact is no longer only a set of weights; it is also the procedure that designs the architecture and training configuration. Classical NAS, however, is typically an external process. A controller proposes architectures, a training loop evaluates them, and engineers deploy the winner. The resulting model does not normally continue to search over its own architecture after deployment. In this respect, AutoML is an ancestor of self-evolution but not always an instance of it.

The analogy becomes tighter when the target of search is an agentic system rather than a neural network. ADAS frames automated design of agentic systems as a search over programs that implement agents \cite{hu2024automated}. The meta-agent writes code for new agent architectures, evaluates them on tasks, and stores successful designs. The search space is Turing-complete because agents are expressed as Python programs. This is NAS for agent software: instead of choosing convolutional cells or transformer blocks, the system searches over prompts, tool calls, memory strategies, decomposition routines, reflection procedures, and control flow. Agent Symbolic Learning makes a related move by treating agents as symbolic networks whose learnable elements are prompts, tools, and pipelines rather than differentiable weights \cite{zhou2024agent}. The method then defines language losses, language gradients, and textual updates analogous to backpropagation.

Evolutionary computation deepens this lineage. NEAT evolves both topology and weights \cite{stanley2002evolving}; novelty search argues that divergent exploration can be more productive than direct objective optimization \cite{lehman2011abandoning}; MAP-Elites and quality-diversity methods maintain archives of high-performing diverse solutions \cite{mouret2015illuminating}; POET co-evolves agents and environments in open-ended processes \cite{wang2019poet}; AI-GAs imagine systems that generate architectures, learning algorithms, and environments \cite{clune2019ai}. Modern self-evolving agents reuse these ideas but replace random or hand-coded mutation operators with LLM-guided semantic variation. The mutation operator now reads code, understands test failures, edits prompts, and proposes structured improvements.

Darwin G{\"o}del Machine is the clearest synthesis. It maintains an archive of agent implementations, samples parent agents, asks a foundation model to modify their code, evaluates the child, and adds improved variants to the archive \cite{zhang2025darwin}. Its reported SWE-bench improvement from 20.0\% to 50.0\% and Polyglot improvement from 14.2\% to 30.7\% indicate that agent architecture can itself be a cumulative evolutionary object. AlphaEvolve provides an industrial-scale counterpart: a Gemini-powered evolutionary coding agent that combines broad exploration with deeper reasoning and MAP-Elites-style quality diversity, reporting improvements in matrix multiplication, data-center scheduling, chip design, and LLM-training kernels \cite{novikov2025alphaevolve}. These systems make explicit what AutoML implied: the design process can become a computational object that is searched, evaluated, and improved.

\subsubsection{The large language model revolution as an enabler of self-evolution}

The rise of large language models changed the practical feasibility of self-evolution. Earlier systems could optimize numeric parameters, discrete architectures, or symbolic rules, but they struggled to operate over the heterogeneous artifacts that make up modern AI systems: prompts, tool APIs, logs, source code, test results, scientific rationales, and natural-language feedback. GPT-4, Claude, Gemini, and related frontier models are not merely stronger predictors of text; they are broadly competent manipulators of software and language-mediated procedures. They can read a stack trace, infer a bug, generate a patch, produce a unit test, critique a solution, rewrite a prompt, summarize a benchmark failure, and propose a new experiment.

This capacity turns language into a universal interface for self-modification. In Self-Refine, the model generates an output, critiques it, and refines it using its own feedback \cite{madaan2023selfrefine}. No weights change, but the output process evolves within an episode. Reflexion adds memory: failed attempts produce verbal reflections that condition later attempts \cite{shinn2023reflexion}. Agent Symbolic Learning pushes the analogy further by defining language gradients over the nodes of an agentic computation graph \cite{zhou2024agent}. RISE trains recursive introspection as a learned multi-turn behavior using reinforcement learning \cite{qu2024rise}. RAGEN decomposes trajectories into state, thinking, actions, and rewards, while identifying the ``Echo Trap'' in which agents appear to improve by repeating themselves rather than reasoning \cite{wang2025ragen}. ReVeal treats self-verification as a first-class learned capability in code agents, optimizing generation and verification together \cite{jin2025reveal}.

The LLM revolution also enabled systems that act on code as an evolutionary substrate. SelfEvolve, an early code-evolution framework, used self-generated API knowledge and execution feedback to iteratively repair code, improving DS-1000 and HumanEval performance \cite{jiang2023selfevolve}. ADAS asks an LLM meta-agent to write new agent code \cite{hu2024automated}. G{\"o}del Agent modifies its own runtime behavior through monkey patches \cite{yin2024godel}. Darwin G{\"o}del Machine and AlphaEvolve use foundation models as semantic mutation operators over agent implementations and algorithmic code \cite{zhang2025darwin,novikov2025alphaevolve}. AI Scientist uses LLM agents to traverse a broader research loop: idea generation, experimental design, code execution, manuscript writing, and peer review \cite{lu2024ai}. These systems would be hard to imagine without models that can manipulate both natural language and executable artifacts.

Importantly, LLMs do not remove the need for evaluation. They make variation cheap and semantically informed, but they also hallucinate, overfit to rubrics, exploit loopholes, and produce plausible but invalid rationales. The strongest systems therefore pair LLM generation with external feedback. Code executors, test suites, benchmark harnesses, environment rewards, and human or model-based judges provide the selection pressure. Absolute Zero's zero-data self-play depends on a code executor that verifies proposed tasks and solutions \cite{zhao2025absolute}. AlphaEvolve depends on automated evaluators \cite{novikov2025alphaevolve}. ReVeal emphasizes reliable self-verification \cite{jin2025reveal}. DGM evaluates modified agents on software-engineering benchmarks \cite{zhang2025darwin}. The lesson is that LLMs make self-evolution possible, but verifiers make it cumulative.

The field's empirical record reflects this verifier dependence. Prompt-only self-feedback can improve many tasks, but it can also stagnate or reinforce errors. Memory-based reflection can increase pass rates when feedback is specific and evaluators are reliable, but memory can grow unbounded. Code-level self-modification can discover novel strategies, but it requires sandboxing, rollback, and careful acceptance criteria. RL-based self-improvement can internalize revision behaviors, but it risks reward hacking, reasoning collapse, or echoing. Open-ended archives can escape local optima, but they are expensive and difficult to govern. These are not isolated problems; they are the core design trade-offs of self-evolving AI.

\subsubsection{From isolated tricks to self-improvement systems}

The most visible early LLM self-improvement methods were local loops. A single answer is critiqued and revised. A failed code sample is debugged. A reflection is stored in memory. These loops are important because they demonstrated that models could use their own outputs as material for improvement. But they are not the end state of the field. Recent work increasingly integrates multiple loops into systems that improve across episodes, across components, or across populations.

A useful progression begins with output-level refinement. Self-Refine improves a candidate output through generate--critique--refine iterations \cite{madaan2023selfrefine}. Reflexion improves future attempts by storing verbal lessons from prior failures \cite{shinn2023reflexion}. SelfEvolve improves code by combining self-generated knowledge with interpreter feedback \cite{jiang2023selfevolve}. These systems modify the context or artifact of a task, not necessarily the agent's source code or weights. They show that feedback can be represented linguistically and reused.

The next step is component-level optimization. Agent Symbolic Learning treats prompts, tools, and pipelines as learnable symbolic weights, allowing a deployed agent to adjust internal components after observing trajectories \cite{zhou2024agent}. DSPy-style prompt compilation, textual backpropagation, and multi-agent prompt evolution belong to the same family. Here the object being improved is not just an answer but the agent's procedure for producing answers. The boundary between inference and learning becomes less clear: the agent can execute tasks while updating symbolic control parameters.

A third step is architecture- and code-level evolution. ADAS uses a meta-agent to write new agents \cite{hu2024automated}. G{\"o}del Agent modifies runtime logic \cite{yin2024godel}. DGM evolves agent implementations through an archive \cite{zhang2025darwin}. AlphaEvolve evolves algorithmic code at scale \cite{novikov2025alphaevolve}. These systems treat software as the genome of intelligence. The genome is readable and editable by a foundation model; evaluation determines which variants survive. This perspective makes agent design cumulative and potentially open-ended.

A fourth step is weight- and curriculum-level self-improvement. RISE trains a model to introspect and revise its reasoning through multi-turn reinforcement learning \cite{qu2024rise}. RAGEN adapts RL to multi-turn agents while warning about echo traps \cite{wang2025ragen}. Absolute Zero removes external training data by letting the model propose tasks and solve them under verifiable rewards \cite{zhao2025absolute}. ReVeal co-optimizes code generation and self-verification \cite{jin2025reveal}. These methods modify neural weights, task curricula, or verification behaviors. They suggest that self-evolution can happen not only at the prompt or code layer but also inside the model's learned policy.

A fifth step is lifecycle-level automation. AI Scientist attempts to automate the scientific research lifecycle, from idea generation to experiment execution and paper writing \cite{lu2024ai}. Production agent frameworks such as AutoGen, LangGraph, Letta/MemGPT, OpenHands, SWE-Agent, CrewAI, MetaGPT, DSPy, OpenEvolve, and related projects vary widely in their degree of self-evolution, but collectively they show how research loops are moving into usable engineering ecosystems. Some frameworks remain mostly static; others support memory evolution, graph modification, programmatic agent control, or evolutionary code search. The open-source ecosystem is therefore both a deployment substrate and a measurement challenge: popularity, stars, and hype do not necessarily correlate with genuine self-evolving capability.

The background thus motivates a new survey category. AI Self-Evolution is not equivalent to LLM prompting, not reducible to AutoML, and not identical to reinforcement learning. It is a cross-layer design pattern in which systems modify themselves through feedback. The rest of this chapter defines that pattern precisely.

\subsection{Core Definition: Formalizing AI Self-Evolution}
\label{subsec:intro-definition}

We define \textbf{AI Self-Evolution} as follows:

\begin{quote}
\textit{An AI system $S$ is self-evolving if it can modify one or more of its own operational components---including prompts, memories, tools, code, architecture, data curriculum, or weights---through feedback-driven loops in order to improve performance on target tasks, subject to an acceptance mechanism that determines whether candidate modifications are retained.}
\end{quote}

This definition emphasizes five elements: a system boundary, mutable components, feedback, a modification operator, and retention. The system boundary specifies what counts as ``self.'' Mutable components specify what may change. Feedback supplies evidence about success or failure. The modification operator proposes changes. Retention decides whether the change becomes part of the future system. Without feedback, change is drift rather than evolution. Without retained modification, the system may adapt within an episode but does not accumulate improvement. Without a system boundary, every external human update would count as self-evolution, making the concept meaningless.

Formally, let an AI system be represented as
\begin{equation}
S_t = (\mathcal{C}_t, \mathcal{M}_t, \mathcal{E}, \mathcal{U}, \mathcal{G}),
\end{equation}
where $\mathcal{C}_t$ denotes mutable components at time $t$, $\mathcal{M}_t$ denotes memory or state, $\mathcal{E}$ denotes an environment or task distribution, $\mathcal{U}$ denotes an objective or utility criterion, and $\mathcal{G}$ denotes governance constraints such as safety policies, sandbox limits, or rollback rules. The mutable component set may include
\begin{equation}
\mathcal{C}_t = \{P_t, R_t, T_t, A_t, K_t, D_t, W_t\},
\end{equation}
where $P_t$ are prompts or instructions, $R_t$ are reasoning routines, $T_t$ are tools and tool descriptions, $A_t$ is agent architecture or control flow, $K_t$ is source code or executable implementation, $D_t$ is data or curriculum, and $W_t$ are neural weights.

During an evolution episode, the system interacts with tasks $x\sim\mathcal{E}$, produces trajectories $\tau_t$, receives feedback $F_t$, proposes a modification $\Delta_t$, and updates its components if an acceptance rule approves:
\begin{align}
\tau_t &= \mathrm{Run}(S_t, x), \\
F_t &= \mathrm{Feedback}(\tau_t, x, \mathcal{U}), \\
\Delta_t &= \mathrm{Modify}(S_t, F_t, \mathcal{M}_t), \\
S_{t+1} &= \mathrm{Accept}(S_t, \Delta_t, F_t, \mathcal{G}).
\end{align}
The modification may be a prompt edit, a memory entry, a tool update, a code patch, an architecture change, a generated training task, or a weight update. The acceptance function may be deterministic, such as keeping a patch only if tests pass, or probabilistic, such as selection in an evolutionary population. It may also include human approval, but the system must contribute materially to proposing or evaluating the change.

A performance-improvement objective can be written as
\begin{equation}
\mathbb{E}_{x\sim\mathcal{E}_{\mathrm{target}}}\left[ U(S_{t+1},x) \right]
\;>\;
\mathbb{E}_{x\sim\mathcal{E}_{\mathrm{target}}}\left[ U(S_t,x) \right] - \epsilon,
\end{equation}
where $\epsilon$ allows for noisy evaluation and exploration. In strict exploitation settings, we require expected improvement. In open-ended settings, however, a system may retain diverse variants that do not immediately dominate the parent but improve coverage, novelty, or future evolvability. For archive-based systems, the objective becomes multi-dimensional:
\begin{equation}
\max_{S\in\mathcal{A}} \left( f(S),\; n(S),\; r(S),\; c(S)^{-1} \right),
\end{equation}
where $\mathcal{A}$ is an archive, $f$ is task fitness, $n$ is novelty or diversity, $r$ is robustness, and $c$ is computational cost. This form captures quality-diversity methods such as MAP-Elites \cite{mouret2015illuminating} and their modern use in AlphaEvolve \cite{novikov2025alphaevolve}.

The phrase ``its own components'' should be interpreted operationally, not metaphysically. A self-evolving agent need not be conscious, autonomous in a legal sense, or free from human-designed constraints. It must merely have a loop by which the system that performs tasks is modified by a process in which that system, or a tightly coupled meta-system, participates. ADAS, for instance, uses a meta-agent to design object-level agents \cite{hu2024automated}. Whether this is one system or two depends on the boundary. We count it as self-evolution when the meta-agent, archive, evaluator, and generated agents form a persistent system whose future designs are conditioned on its own search history. Similarly, Reflexion's actor, evaluator, and reflection model can be viewed as a composite system whose memory updates alter future behavior \cite{shinn2023reflexion}.

This boundary-based view avoids two extremes. On one extreme, every training procedure could be called self-evolution because it improves a model. That is too broad. Standard supervised learning updates weights by an externally specified optimizer over a fixed dataset; the model does not propose changes to its own learning process or components. On the other extreme, only systems that rewrite their own source code would qualify. That is too narrow. Modern AI behavior is often controlled by prompts, retrieval state, memories, tools, and workflow graphs. If those components are modified by feedback-driven loops, the system has evolved at the level that matters for behavior.

We therefore distinguish levels of self-evolution:

\begin{enumerate}
    \item \textbf{Output-level evolution}: the system revises a candidate output within a task episode, as in generate--critique--refine loops.
    \item \textbf{Memory-level evolution}: the system stores feedback, reflections, or experiences that alter future behavior.
    \item \textbf{Prompt/tool-level evolution}: the system edits prompts, tool descriptions, or symbolic parameters.
    \item \textbf{Architecture-level evolution}: the system changes control flow, modular composition, or agent topology.
    \item \textbf{Code-level evolution}: the system writes or patches executable source code implementing itself or its descendants.
    \item \textbf{Data/curriculum-level evolution}: the system generates or selects tasks that shape future learning.
    \item \textbf{Weight-level evolution}: the system updates neural parameters using feedback from self-generated or self-mediated experience.
    \item \textbf{Population-level evolution}: the system maintains multiple variants and selects among them over time.
\end{enumerate}

These levels are not mutually exclusive. ReVeal combines code generation, self-verification, tool feedback, and reinforcement learning \cite{jin2025reveal}. Absolute Zero combines task generation, solving, executable verification, and weight updates \cite{zhao2025absolute}. DGM combines code-level modification with population-level archiving \cite{zhang2025darwin}. AI Scientist combines idea generation, experiment code, paper writing, and review loops \cite{lu2024ai}. A mature self-evolving system will likely contain several levels simultaneously.

Feedback is the central resource. It can be internal, external, symbolic, scalar, textual, executable, or social. Internal feedback includes self-critique and introspection. External feedback includes unit tests, environment rewards, benchmark scores, human judgments, peer-review scores, and market signals. Symbolic feedback includes language gradients and reflection notes. Scalar feedback includes rewards, pass/fail signals, and performance metrics. Executable feedback includes stack traces, assertions, and verifier outputs. Social feedback includes user corrections and community adoption. The reliability of self-evolution depends heavily on the reliability and alignment of this feedback. A weak evaluator produces weak evolution; a manipulable evaluator selects for manipulation.

The modification operator is equally important. In classical evolution, mutation may be random; in gradient descent, the update direction is computed from differentiable loss; in AutoML, search operators are hand-designed or controller-generated. In modern self-evolving AI, LLMs often implement the modification operator:
\begin{equation}
\Delta_t = \mathrm{LLM}_{\phi}(\mathrm{serialize}(S_t), F_t, H_t),
\end{equation}
where $H_t$ is history, archive context, or instructions. This operator is powerful because it can make semantically coherent edits across heterogeneous artifacts. It is dangerous because it is stochastic, prompt-sensitive, and capable of introducing hidden failures. The same LLM that patches a bug can remove a guardrail; the same agent that improves a benchmark can overfit to its public tests.

Retention mechanisms define whether improvements accumulate. A naive system may overwrite itself after every self-critique, producing drift. A more robust system evaluates candidate changes on validation tasks, compares against baselines, stores rollback points, and preserves diversity. Reflexion's memory buffer retains reflections but must manage context limits \cite{shinn2023reflexion}. G{\"o}del Agent accepts modifications only after validation and can roll back failed patches \cite{yin2024godel}. DGM stores variants in an archive rather than a single mutable lineage \cite{zhang2025darwin}. AlphaEvolve uses quality-diversity selection to retain elites in behavioral niches \cite{novikov2025alphaevolve}. Retention is where learning becomes evolution.

Finally, self-evolution must be evaluated over time. A static benchmark score after one update is insufficient. We must ask whether the system improves monotonically or cyclically, whether gains transfer across tasks, whether it maintains safety constraints, whether it collapses into repetitive behavior, whether it overfits to evaluators, and whether the cost of evolution is justified by the benefit. For an evolution trace $S_0,S_1,\ldots,S_T$, relevant metrics include cumulative improvement,
\begin{equation}
\Delta_T = U(S_T)-U(S_0),
\end{equation}
per-step efficiency,
\begin{equation}
\eta_T = \frac{U(S_T)-U(S_0)}{\sum_{t=0}^{T-1} c_t},
\end{equation}
and robustness under held-out distributions,
\begin{equation}
\rho_T = \mathbb{E}_{x\sim\mathcal{E}_{\mathrm{heldout}}}[U(S_T,x)] - \mathbb{E}_{x\sim\mathcal{E}_{\mathrm{train}}}[U(S_T,x)].
\end{equation}
These metrics foreshadow the evaluation chapter of the survey, where we argue that self-evolving systems require benchmarks that measure trajectories, not just endpoints.

\subsection{Related Concepts and Boundaries}
\label{subsec:intro-related}

AI Self-Evolution overlaps with several established fields, but it is not identical to any of them. This subsection draws boundaries because the term ``self-improvement'' is often used loosely. A system may improve after deployment, learn online, adapt to a user's preferences, tune hyperparameters, or train on self-supervised losses without satisfying the stronger criteria of self-evolution. Conversely, a system may be self-evolving even if its neural weights remain fixed, provided that feedback-driven changes to prompts, memory, tools, code, or architecture alter future behavior.

\paragraph{Continual learning.} Continual learning studies models that learn from a sequence of tasks or data distributions without catastrophic forgetting. Its central questions include stability-plasticity trade-offs, replay, regularization, task inference, and lifelong adaptation. A continual learner updates parameters over time, but the update rule is usually fixed by the designer. The model does not necessarily inspect failures, propose new tools, rewrite its architecture, or choose its own curriculum. Continual learning becomes self-evolution when the system participates in modifying its own learning machinery or task interface, not merely when more data are streamed through a fixed optimizer.

\paragraph{Online learning.} Online learning updates predictions as examples arrive. It is often formalized with regret bounds against a comparator. The learner receives an input, predicts, observes loss, and updates. This is feedback-driven, but the feedback is typically used by a fixed algorithm such as stochastic gradient descent, multiplicative weights, or bandit updates. Self-evolution requires an additional layer: the system may change the representation, prompt, code, architecture, or update rule itself. RISE and RAGEN are closer to self-evolution than generic online learning because they learn multi-turn introspection or agent policies, but even they must be analyzed by whether the learned behavior changes future self-modification capacity \cite{qu2024rise,wang2025ragen}.

\paragraph{Self-supervised learning.} Self-supervised learning constructs labels from data itself, as in masked language modeling, contrastive learning, or next-token prediction. It is ``self'' in the sense that supervision is derived automatically, not manually annotated. It is not necessarily self-evolving because the model does not modify its own operational components through feedback about task performance. Absolute Zero is an instructive contrast: it generates tasks and solutions, verifies them with a code executor, and updates via reinforcement learning \cite{zhao2025absolute}. This is more than self-supervision because the model shapes its own curriculum under verifiable rewards.

\paragraph{Meta-learning.} Meta-learning trains systems that learn how to learn across tasks. It is closely related to self-evolution because it concerns adaptation mechanisms. However, much meta-learning remains offline: a meta-learner is trained on task distributions and then adapts using a fixed inner-loop procedure. Self-evolution emphasizes post-deployment or within-system modification of components, often including symbolic and code artifacts not captured by differentiable meta-learning. ADAS can be interpreted as meta-learning over agent designs, but it is also self-evolutionary because the search history and generated agents condition future designs \cite{hu2024automated}.

\paragraph{AutoML and NAS.} AutoML automates model design, training, and selection. NAS searches architecture spaces. These fields are direct predecessors, but most AutoML pipelines are external optimizers. A human invokes the optimizer, waits for a candidate model, and deploys it. AI Self-Evolution internalizes part of that optimizer into the system lifecycle. When a meta-agent designs new agents, when an archive of agent implementations accumulates, or when a coding system mutates its own algorithms under evaluators, AutoML becomes agentic and self-referential.

\paragraph{Reinforcement learning.} Reinforcement learning optimizes policies through reward signals. It is foundational for self-evolution, especially in RISE, RAGEN, Absolute Zero, and ReVeal \cite{qu2024rise,wang2025ragen,zhao2025absolute,jin2025reveal}. Yet RL alone does not imply self-evolution. A policy trained by PPO on environment rewards may become more competent, but it does not necessarily modify its prompts, tools, code, architecture, or curriculum. Self-evolution appears when the policy controls or participates in the modification process, such as proposing tasks, generating tests, revising reasoning, or altering agent software.

\paragraph{RLHF and preference optimization.} Reinforcement learning from human feedback aligns models to preferences \cite{christiano2017deep,ziegler2019fine}. It uses feedback, but the process is usually externally orchestrated: humans label data, engineers train reward models, and training pipelines update weights. Constitutional AI and self-critique methods move closer to self-evolution by using model-generated critiques and revisions under a constitution \cite{bai2022constitutional}. Still, the distinction depends on whether the deployed system can continue modifying its own components under feedback rather than merely being trained once with synthetic preferences.

\paragraph{Agent memory and personalization.} Many agent frameworks store memories, user preferences, or conversation summaries. Memory can be self-evolutionary if it is generated from feedback and retained to improve future behavior. Letta/MemGPT-like memory systems are therefore partial self-evolution mechanisms. But memory alone is not enough. A logging system that records transcripts without using them for future modification is not self-evolving. A recommender that stores user preferences may be adaptive, but if its update rule and task interface are fixed, it occupies a narrower category.

\paragraph{Program synthesis and automated repair.} Program synthesis generates code from specifications; automated program repair patches bugs according to tests. Modern LLM code agents blur this boundary. SelfEvolve, ReVeal, AlphaEvolve, and DGM all use program synthesis or repair as part of self-evolution \cite{jiang2023selfevolve,jin2025reveal,novikov2025alphaevolve,zhang2025darwin}. The distinguishing feature is recursive use: code generation is not only used to solve an external programming task but also to improve the system's own future problem-solving procedure.

\begin{table}[t]
\centering
{\TableTight
\begin{tabularx}{\textwidth}{@{}L{0.16\textwidth}L{0.20\textwidth}L{0.23\textwidth}Y@{}}
\toprule
\textbf{Concept} & \textbf{Primary object of change} & \textbf{Typical feedback} & \textbf{Difference from AI Self-Evolution} \\
\midrule
Continual learning & Neural parameters over task streams & New data, task labels, replay losses & Usually uses a fixed update rule; may not modify prompts, tools, code, architecture, or curriculum. \\
Online learning & Predictor state after each example & Immediate loss or reward & Focuses on regret under fixed algorithms; self-evolution can change the algorithmic machinery itself. \\
Self-supervised learning & Representations and weights & Automatically generated labels & Supervision is automatic, but the system does not necessarily evaluate and modify its own operational components. \\
Meta-learning & Adaptation procedure or initialization & Distribution of training tasks & Often learned offline; self-evolution emphasizes runtime feedback loops and retained system modifications. \\
AutoML/NAS & Model pipelines or architectures & Validation performance & Usually an external design optimizer; becomes self-evolutionary when integrated into the agent/system lifecycle. \\
Reinforcement learning & Policy parameters & Environment rewards & RL is an optimization substrate; self-evolution additionally concerns self-modification of components and procedures. \\
RLHF/preference tuning & Model behavior and alignment & Human or AI preference labels & Usually externally orchestrated; self-evolution requires ongoing system-mediated modification. \\
Agent memory & Episodic or semantic context & User interactions, reflections & Partial self-evolution if memory is feedback-derived and behaviorally active; otherwise only logging/personalization. \\
Program synthesis/repair & Task-specific code & Specifications, tests, errors & Self-evolution when generated or repaired code changes the future agent/system, not just a one-off output. \\
\bottomrule
\end{tabularx}}
\caption{AI Self-Evolution in relation to neighboring concepts. The boundaries are porous: many modern systems combine several rows, but self-evolution specifically requires feedback-driven modification of the system's own operational components.}
\label{tab:related-concepts}
\end{table}

The comparison in Table~\ref{tab:related-concepts} suggests a general rule. If the system merely changes its output, it is adaptation. If it changes an internal state that affects future outputs, it may be learning. If it changes the components or procedures by which it learns, reasons, acts, or designs successors, under feedback and retention, it is self-evolution. This rule captures both lightweight prompt/memory methods and heavyweight code/weight/population methods.

Several borderline cases deserve attention. First, an LLM using chain-of-thought prompting does not self-evolve if each response is independent. If the chain is critiqued, revised, stored, and reused, it may become output- or memory-level self-evolution. Second, a human editing an agent's prompt after observing a failure is not AI self-evolution, though it may be part of a human-in-the-loop evolution system. If the agent proposes the edit and the human only approves it, the system is partially self-evolving. Third, a benchmark-tuned model release is not self-evolving after deployment, but the training organization may implement an external evolution pipeline. Our survey focuses on systems whose loop is technically represented and can be analyzed as part of the AI system.

Another boundary concerns autonomy. A self-evolving system need not be fully autonomous. Human approval, safety filters, sandboxing, and governance are often desirable. In fact, unrestricted autonomy is a safety risk. We therefore define self-evolution by participation in modification, not by absence of human oversight. A system that drafts code patches for itself, evaluates them in a sandbox, and requests approval before deployment still exhibits self-evolutionary structure. The acceptance function includes humans as part of governance $\mathcal{G}$.

A final boundary concerns improvement. Evolutionary systems can regress. A single failed modification does not disqualify a system, just as a biological lineage can accumulate deleterious mutations. What matters is that the system contains a feedback-driven selection process intended to improve performance or diversity over time. Evaluation should measure the trajectory, including failures, rather than assume monotonic progress.

\subsection{Contributions of This Survey}
\label{subsec:intro-contributions}

This survey makes four main contributions. First, we introduce a \textbf{Five Evolution Loops} framework that organizes the field by the mechanism through which feedback becomes retained change. Second, we analyze more than nineteen research and open-source systems, separating genuine self-evolution mechanisms from static agent orchestration and hype-driven adoption. Third, we unify work across AutoML/NAS, LLM self-improvement, evolutionary computation, reinforcement learning, program synthesis, and agent frameworks under a common formal abstraction. Fourth, we map the researcher and institution network that produced the field, highlighting the convergence of the evolutionary-computation lineage, the language-agent reasoning lineage, and the LLM self-improvement lineage.

\subsubsection{Contribution 1: The Five Evolution Loops framework}

The first contribution is a framework that classifies self-evolving AI by the loop that performs modification and selection. We use five loops because they capture the dominant mechanisms in the literature while remaining general enough to compare hybrid systems.

\paragraph{Loop I: Reflexive output and memory evolution.} The system improves by critiquing outputs, storing feedback, and conditioning future attempts. Self-Refine and Reflexion are representative \cite{madaan2023selfrefine,shinn2023reflexion}. The modified component is usually the current output, an episodic memory, or a reflection buffer. The evaluator may be the same model, an external reward model, tests, or an environment. This loop is low-cost and easy to deploy, but it is limited by context length, self-critique quality, and evaluator reliability.

\paragraph{Loop II: Symbolic component evolution.} The system treats prompts, tools, workflows, and pipelines as learnable symbolic parameters. Agent Symbolic Learning is the central example: it defines language losses, language gradients, and optimizers for prompt, tool, and pipeline nodes \cite{zhou2024agent}. DSPy-like prompt compilation and textual backpropagation also fit here. This loop bridges gradient-based learning and agent engineering. It makes the agent's design surface optimizable without requiring neural weight updates.

\paragraph{Loop III: Verification-driven code evolution.} The system generates, tests, repairs, and retains code. SelfEvolve uses self-generated knowledge plus interpreter errors \cite{jiang2023selfevolve}; ReVeal co-optimizes generation and self-verification \cite{jin2025reveal}; AlphaEvolve uses automated evaluators to select algorithmic programs \cite{novikov2025alphaevolve}. This loop is powerful because code can be mechanically checked. Its limitation is domain dependence: reliable verifiers are easier for programming, mathematics, and simulation than for open-ended social tasks.

\paragraph{Loop IV: Architecture and agent-design evolution.} The system searches over agent architectures, control flow, tool use, and self-modifying code. ADAS, G{\"o}del Agent, and DGM are representative \cite{hu2024automated,yin2024godel,zhang2025darwin}. The modified component is the agent itself. This loop is closest to the intuitive idea of self-redesign. It can discover strategies not anticipated by humans, but it raises cost, safety, and reproducibility challenges.

\paragraph{Loop V: Curriculum, weight, and population evolution.} The system changes the data distribution, training curriculum, neural weights, or population archive. RISE, RAGEN, Absolute Zero, DGM, and AlphaEvolve all instantiate parts of this loop \cite{qu2024rise,wang2025ragen,zhao2025absolute,zhang2025darwin,novikov2025alphaevolve}. The loop can produce deeper behavioral changes than prompt-only methods, but it requires careful reward design and protection against collapse. Absolute Zero's dual role of task proposer and solver is especially important because it removes the assumption of human-curated training data.

The loops can be summarized as
\begin{equation}
\mathrm{SelfEvolution} = \mathrm{Observe} \rightarrow \mathrm{Interpret} \rightarrow \mathrm{Modify} \rightarrow \mathrm{Verify} \rightarrow \mathrm{Retain},
\end{equation}
with each loop differing in the representation modified and the verifier used. The framework is not a taxonomy of papers; it is a taxonomy of mechanisms. A single paper may instantiate multiple loops. For example, DGM combines architecture/code evolution with population archives; ReVeal combines verification-driven code evolution with weight-level RL; AI Scientist combines idea, code, manuscript, and review loops.

\begin{table}[t]
\centering
{\TableTight
\begin{tabularx}{\textwidth}{@{}L{0.18\textwidth}L{0.19\textwidth}L{0.25\textwidth}Y@{}}
\toprule
\textbf{Loop} & \textbf{Mutable component} & \textbf{Representative systems} & \textbf{Typical failure modes} \\
 \midrule
Reflexive output and memory & Outputs, reflections, memory buffers & Self-Refine, Reflexion, ReAct-style agents & Generic critiques, memory bloat, context overfitting, local optima. \\
Symbolic component & Prompts, tools, pipelines, workflow nodes & Agent Symbolic Learning, DSPy, EvoMAC-like textual optimization & Prompt overfitting, unstable language gradients, high LLM-call cost. \\
Verification-driven code & Candidate code, tests, repair patches & SelfEvolve, ReVeal, AlphaEvolve, code agents & Test overfitting, sandbox escape risk, verifier incompleteness. \\
Architecture and agent design & Agent code, control flow, tool topology & ADAS, G{\"o}del Agent, DGM & Unsafe self-modification, irreproducible search, expensive evaluation. \\
Curriculum, weight, and population & Training tasks, policy weights, archives & RISE, RAGEN, Absolute Zero, DGM, AlphaEvolve & Reward hacking, echo trap, mode collapse, compute cost. \\
\bottomrule
\end{tabularx}}
\caption{The Five Evolution Loops framework used throughout this survey. The loops classify mechanisms rather than mutually exclusive systems.}
\label{tab:five-loops}
\end{table}

\subsubsection{Contribution 2: Analysis of research systems and open-source projects}

The second contribution is an evidence-oriented analysis of more than nineteen systems. The research core includes Self-Refine, Reflexion, Agent Symbolic Learning, ADAS, G{\"o}del Agent, Darwin G{\"o}del Machine, SelfEvolve, RISE, RAGEN, Absolute Zero, ReVeal, AlphaEvolve, FunSearch, AI Scientist, STaR, Constitutional AI, Tree of Thoughts, ReAct, SWE-agent \cite{sweagent2024}, and related software-engineering agents. The open-source and framework ecosystem includes DSPy, AutoGen, LangGraph, Letta/MemGPT, CrewAI, MetaGPT, OpenHands, OpenEvolve, AutoGPT, smolagents, and project-specific implementations of the papers above.

Our analysis emphasizes mechanism rather than popularity. The star-analysis report used for this survey shows that GitHub stars can be a poor proxy for technical quality. AutoGPT accumulated massive attention, but its star-quality score was estimated far below smaller research-oriented projects. OpenEvolve, DSPy, SWE-Agent, OpenHands, LangGraph, AutoGen, CrewAI, and MetaGPT exhibit different propagation patterns: academic, steady, hype-driven, or community-driven. This matters because self-evolution is a technical property, not a marketing label. A project may be popular because it is easy to demo, while a quieter project may contain the more important self-evolution mechanism.

We therefore examine each system along dimensions such as mutable component, feedback source, retention mechanism, evaluation benchmark, openness, safety controls, and transfer. Reflexion is strong evidence for memory-level verbal feedback because it reports large HumanEval and interactive-task gains under external evaluation \cite{shinn2023reflexion}. Self-Refine demonstrates that same-model critique can improve diverse tasks by roughly 20\% on average, but it also illustrates the limits of self-feedback without external verification \cite{madaan2023selfrefine}. Agent Symbolic Learning reports improvements on HotPotQA, MATH, HumanEval, software-development executability, and creative writing by optimizing symbolic agent components \cite{zhou2024agent}. DGM reports striking software-engineering improvements through open-ended code-level evolution \cite{zhang2025darwin}. AlphaEvolve reports algorithmic and infrastructure improvements under automated evaluation \cite{novikov2025alphaevolve}. Absolute Zero demonstrates that self-generated curricula can support zero-data reasoning improvements \cite{zhao2025absolute}.

The open-source ecosystem is equally important for understanding adoption. AutoGen supports multi-agent conversation and human-in-the-loop workflows; LangGraph offers state graphs that can be modified or orchestrated programmatically; Letta/MemGPT emphasizes long-term memory; DSPy compiles prompts and demonstrations; SWE-Agent and OpenHands operationalize software-engineering agents; OpenEvolve adapts evolutionary coding ideas for the community. CrewAI and MetaGPT popularize role-based agent orchestration, but many deployments remain static unless combined with feedback-driven modification. Distinguishing static orchestration from self-evolution is one of the survey's practical aims.

\subsubsection{Contribution 3: Cross-domain unification}

The third contribution is a cross-domain unification. The same abstract loop appears under different names in different fields. In AutoML it is architecture search. In evolutionary computation it is variation and selection. In LLM prompting it is critique and refinement. In reinforcement learning it is policy improvement. In program synthesis it is generate-and-test. In software engineering it is automated repair. In agent frameworks it is memory and workflow adaptation. In AI-for-science it is hypothesis generation, experimentation, and review.

We formalize these domains using a shared tuple: representation, variation operator, feedback channel, retention rule, and governance constraint. For NAS, the representation is a neural architecture, variation is architecture mutation or controller sampling, feedback is validation accuracy, retention is selection, and governance is compute budget. For Reflexion, the representation is memory, variation is verbal reflection, feedback is evaluator reward, retention is memory update, and governance is context management. For DGM, the representation is agent source code, variation is LLM patching, feedback is benchmark performance, retention is archive insertion, and governance is sandboxing and selection. For Absolute Zero, the representation includes task curriculum and weights, variation is self-generated tasks and solutions, feedback is code execution, retention is RL update, and governance is verifier design.

This unification reveals shared bottlenecks. Search cost appears in NAS, ADAS, DGM, and AlphaEvolve. Evaluator reliability appears in Self-Refine, Reflexion, ReVeal, AI Scientist, and RLHF. Diversity maintenance appears in evolutionary computation, DGM, AlphaEvolve, and open-ended science. Reward hacking appears in RL, code generation, and benchmark-driven agent design. Safety constraints appear wherever code or tools can be modified. The terminology differs, but the engineering problems rhyme.

The unification also reveals transfer opportunities. NAS developed weight sharing and benchmark suites such as NAS-Bench; self-evolving agents need analogous reusable evaluation landscapes. Quality-diversity algorithms developed archives and novelty metrics; code-evolving agents can use them to avoid local optima. Program repair developed test-generation and patch-validation methods; self-modifying agents need them for safe code evolution. RL developed off-policy evaluation and reward-model diagnostics; self-play and self-verification methods need them to detect collapse. Software engineering developed CI/CD, sandboxing, and rollback; self-evolving AI needs these as governance primitives.

\subsubsection{Contribution 4: Researcher network mapping}

The fourth contribution is a map of the researcher network. The field is young, but the sources reveal three converging lineages. The evolutionary-computation lineage runs from Risto Miikkulainen, Kenneth Stanley, Joel Lehman, and Jeff Clune through NEAT, novelty search, quality diversity, POET, AI-GAs, ADAS, DGM, and AlphaEvolve-related work \cite{stanley2002evolving,lehman2011abandoning,clune2019ai,hu2024automated,zhang2025darwin}. The language-agent reasoning lineage runs through Karthik Narasimhan, Shunyu Yao, Noah Shinn, ReAct, Tree of Thoughts, Reflexion, SWE-bench, and SWE-agent \cite{yao2023react,yao2023tree,shinn2023reflexion}. The LLM self-improvement lineage includes Aman Madaan, Peter Clark, Sean Welleck, Eric Zelikman, and work on Self-Refine, STaR, self-rewarding models \cite{selfrewarding2024}, and constitutional self-critique \cite{madaan2023selfrefine,zelikman2022star,bai2022constitutional}.

These lineages now overlap institutionally and methodologically. Jeff Clune's group appears in ADAS, DGM, AI Scientist, and open-ended evolutionary AI. Shengran Hu and Cong Lu connect ADAS and DGM; Cong Lu also appears in AI Scientist-related work. Shunyu Yao and Karthik Narasimhan connect ReAct, Reflexion, Tree of Thoughts, SWE-agent, and broader language-agent evaluation. AIWaves and Zhejiang University contribute Agent Symbolic Learning, which explicitly bridges agent frameworks and textual optimization. Google DeepMind contributes FunSearch and AlphaEvolve, scaling LLM-guided evolution to mathematics and infrastructure. Microsoft Research contributes ReVeal's verification-centered code-agent RL. Tsinghua and LeapLabTHU contribute Absolute Zero's zero-data self-play reasoning.

Mapping the network is not merely sociological. It explains why ideas appear together. DGM's archive design makes sense against Clune's open-endedness background. Reflexion's memory-based verbal RL fits the Princeton/Northeastern language-agent lineage. Agent Symbolic Learning reflects industry-academic interest in deployable agent frameworks. AlphaEvolve reflects DeepMind's AI-for-science and systems-optimization agenda. AI Scientist reflects Sakana AI's nature-inspired, open-ended automation focus. Understanding these lineages helps researchers locate missing bridges, such as stronger connections between formal program repair and agent self-modification, or between continual-learning theory and prompt/tool evolution.


\subsubsection{Running examples and evidence used throughout the survey}

To make the survey concrete, we use several running examples that span the five loops. These examples are not chosen merely because they are well known; they are chosen because each exposes a different design decision in self-evolving AI. Self-Refine demonstrates the minimal loop: no external training, no explicit memory across tasks, and no code modification, only iterative critique and revision inside a task episode \cite{madaan2023selfrefine}. Reflexion adds an explicit memory buffer and an evaluator, turning feedback from one attempt into context for the next \cite{shinn2023reflexion}. Agent Symbolic Learning moves from episodic improvement to component optimization by representing an agent as a symbolic network with prompts, tools, and pipeline nodes \cite{zhou2024agent}. ADAS then makes the architecture itself the search object, asking a meta-agent to write new agent programs \cite{hu2024automated}. DGM extends that search into an open-ended archive of self-improving agent implementations \cite{zhang2025darwin}. AlphaEvolve scales the same broad idea to industrial algorithm discovery and systems optimization \cite{novikov2025alphaevolve}. Absolute Zero and ReVeal show how self-evolution can also alter the training distribution and policy weights through verifiable rewards \cite{zhao2025absolute,jin2025reveal}.

These examples reveal that self-evolution is not a single point on a capability ladder. A prompt-level loop may be more reliable than a code-level loop in settings where feedback is subjective and safety constraints are tight. A code-level loop may be more powerful in programming tasks because unit tests provide crisp feedback. A weight-level loop may produce deeper generalization, but it is slower, more expensive, and more vulnerable to reward misspecification. An archive-based loop may be the best way to avoid local optima, yet it complicates deployment because the system no longer has one canonical state. The survey therefore avoids ranking methods by apparent sophistication alone. Instead, it asks which loop is appropriate for which task, verifier, cost regime, and governance requirement.

Consider Self-Refine as a first running example. Its generate--critique--refine structure is simple enough to be implemented with a single model and a prompt template. The method reported improvements across code optimization, mathematical reasoning, sentiment reversal, dialogue, poetry, readability, and constrained generation \cite{madaan2023selfrefine}. Its importance lies less in any one benchmark number than in the demonstration that a foundation model can use its own natural-language feedback as an optimization signal. But Self-Refine also exposes a core weakness: the same model that produced an error may fail to diagnose it. Without an external evaluator, self-critique can become rationalization. In our framework, Self-Refine is therefore a high-accessibility, low-assurance instance of Loop I.

Reflexion strengthens the loop by separating actor, evaluator, and self-reflection roles \cite{shinn2023reflexion}. The actor attempts a task; the evaluator assigns success or failure; the reflection model writes a verbal lesson; memory carries that lesson forward. On code and interactive tasks, this architecture can produce large gains because failures are turned into reusable advice. Reflexion's reported HumanEval result, along with improvements on AlfWorld and WebShop, made verbal reinforcement learning a central paradigm. Yet Reflexion also makes visible the memory-management problem. If every failure becomes a memory, context grows; if memories are summarized too aggressively, useful detail is lost; if reflections are generic, they add tokens without adding guidance. Later systems inherit this trade-off whenever they store textual feedback.

Agent Symbolic Learning is our main running example for Loop II because it changes the object of optimization from answers or memories to the symbolic components of an agent \cite{zhou2024agent}. The method's language loss and language gradients are analogies to differentiable learning, but the analogy is operationally useful. A failing trajectory can be converted into a textual diagnosis; the diagnosis can be backpropagated through prompt, tool, and pipeline nodes; each node can then be updated by an LLM optimizer. The reported gains on HotPotQA, MATH, HumanEval, software-development executability, and creative writing suggest that agent components can be optimized in a structured way. At the same time, the method raises new questions: What is a learning rate in language space? How stable are language gradients? How should conflicting textual gradients be aggregated? How do we prevent a prompt optimizer from overfitting to a small development set? These questions connect symbolic agent optimization to the older theory of optimization, but they do not reduce to it.

ADAS is our running example for automated agent-design search \cite{hu2024automated}. It treats Python programs implementing agents as the search space. This matters because the space is not just large; it is expressive enough to encode arbitrary computation. A meta-agent can invent new decomposition procedures, tool-use policies, memory layouts, voting schemes, verification routines, or recursive calls. The promise is clear: human designers no longer need to enumerate all plausible agent architectures. The risk is equally clear: generated agents may be brittle, unsafe, or hard to interpret. ADAS reported that automatically discovered agents can outperform hand-designed baselines and transfer across domains or model backbones. The survey uses ADAS to discuss when search over agent code is better than manual engineering and when the evaluation cost becomes prohibitive.

G{\"o}del Agent is a running example for runtime self-reference \cite{yin2024godel}. Instead of maintaining a population of separate descendants, it lets an agent inspect and modify its own logic through monkey patching. This design is close to the intuitive notion of a self-rewriting agent. It also brings software-engineering hazards to the foreground. Runtime patching can introduce nondeterminism, hidden state dependencies, portability problems, and rollback complexity. The method uses validation and rollback, but the general lesson is broader: as self-evolution moves from prompts to executable code, the system must adopt practices from secure software deployment. Patches need provenance, tests, sandboxing, audit logs, and clear acceptance criteria.

DGM is our running example for open-ended agent evolution \cite{zhang2025darwin}. Its archive converts self-improvement from a single lineage into a population process. This is crucial because one agent variant may improve tool use while another improves planning, and a third may improve error recovery. A single greedy lineage could discard variants that are not immediately best but contain useful building blocks. The archive preserves diversity and supports cumulative innovation. DGM's reported improvement on SWE-bench and Polyglot makes it an important empirical milestone. The survey uses DGM to analyze archive sampling, diversity bonuses, cross-domain transfer, and the difference between improving an agent's score and improving the evolvability of the agent population.

AlphaEvolve is our running example for quality-diversity search with strong industrial evaluation \cite{novikov2025alphaevolve}. It combines multiple Gemini models, automated evaluators, and MAP-Elites-style retention. Its reported results on matrix multiplication, data-center scheduling, chip design, and attention kernels show why code and algorithms are attractive domains for self-evolution: candidates can be executed, measured, compared, and retained. Yet AlphaEvolve also illustrates the infrastructural barrier. Many domains lack reliable, cheap, automated evaluators. A system that can improve a kernel by 23\% under a precise performance benchmark may not be able to improve a policy argument, a medical decision, or a social interaction with the same confidence. The availability of verifiers is thus a central axis of the field.

Absolute Zero is our running example for self-generated curricula \cite{zhao2025absolute}. The model proposes tasks, solves them, and receives verification from a code executor. This dual role resembles self-play in games, but with programming and reasoning tasks replacing board positions. The paradigm is important because human-curated data may become a bottleneck for further improvement. If a model can generate challenging, solvable, diverse tasks, it can create its own training signal. But the same setup risks triviality, unsolvable tasks, mode collapse, or curriculum myopia. The proposer and solver can collude in a degenerate region of task space unless rewards encourage both difficulty and solvability. Absolute Zero therefore highlights the governance problem inside the training loop itself.

ReVeal is our running example for reliable self-verification in code agents \cite{jin2025reveal}. It treats verification not as a post-hoc tool call but as a capability that should be optimized along with generation. This is an important shift. Many agents can write code; fewer can design effective tests, interpret failures, and decide when to stop. ReVeal's turn-level credit assignment and generation-verification co-evolution point toward agents that learn how to use tools for self-correction. The broader lesson is that self-evolution depends not only on better generators but also on better critics, verifiers, and credit assignment.

AI Scientist is a running example for lifecycle-level self-evolution \cite{lu2024ai}. It attempts to automate a full research cycle: idea generation, novelty checking, experiment design, code execution, manuscript writing, and review. It broadens the survey beyond coding agents and prompt optimizers. In scientific discovery, the object being improved may be a hypothesis, experiment, paper, or research agenda. The feedback may be experimental results, novelty scores, or peer-review rubrics. AI Scientist also demonstrates the limits of current systems: generated papers may lack theoretical depth, reviewers may be unreliable, and domains may be narrow. These limitations are not incidental; they show that long-horizon self-evolution requires reliable intermediate evaluations at every stage of a complex workflow.

Together, these running examples create a ladder from local revision to open-ended system redesign. However, the ladder should not be read as a chronology in which later methods simply replace earlier ones. In practice, strong systems will compose loops. A DGM-style agent may use Reflexion-like memory during evaluation. An ADAS-generated architecture may include Self-Refine subroutines. A ReVeal-trained code model may be embedded inside AlphaEvolve-like population search. An AI Scientist may use Agent Symbolic Learning to optimize its review prompts and DGM-like archives to preserve successful experimental strategies. The survey's central claim is that the future field will be compositional: self-evolution will be an architecture of interacting loops.

\subsubsection{Why the field needs a survey now}

The need for a survey arises from conceptual fragmentation. The same mechanism is being rediscovered under different labels: self-refinement, verbal reinforcement learning, recursive introspection, agent symbolic learning, automated design of agentic systems, code evolution, self-play reasoning, reliable self-verification, open-ended discovery, and AI scientist automation. Without a common vocabulary, it is difficult to compare claims. A paper may report that an agent ``self-improves'' when it only revises outputs within a single context window. Another may claim ``self-evolution'' because it trains on self-generated data, even though the deployed system cannot modify its own components. A third may evolve code in a population but not address safety or governance. The survey provides a vocabulary for saying precisely what changed, who changed it, using which feedback, retained by which rule, and evaluated on which distribution.

The field also needs a survey because empirical results are difficult to compare. Self-Refine reports average improvements across diverse tasks; Reflexion reports code and interactive-task gains; Agent Symbolic Learning reports improvements across QA, math, code, software development, and creative writing; DGM reports SWE-bench and Polyglot improvements; AlphaEvolve reports scientific and infrastructure discoveries; Absolute Zero reports zero-data reasoning gains; ReVeal reports multi-turn code-agent improvements. These results do not share a benchmark suite, cost accounting method, safety protocol, or standard definition of an evolution step. A static model leaderboard cannot capture the difference between a one-time prompt refinement and a 1,000-iteration archive search. Survey-level synthesis is necessary before the community can build fair benchmarks.

A third reason is the gap between research and open-source adoption. Popular agent frameworks often advertise autonomy, memory, or multi-agent collaboration, but many deployments remain static. A role-based CrewAI or MetaGPT workflow may coordinate agents without modifying itself. AutoGen conversations may adapt through dialogue but not necessarily retain structural changes. LangGraph can represent dynamic workflows, yet the graph usually changes because a developer edits it. Letta/MemGPT evolves memory, but not necessarily tools or code. DSPy compiles prompts, but its self-evolution depends on whether compilation is integrated into a feedback loop. OpenEvolve and related projects bring evolutionary coding to practitioners, but verifier quality and benchmark choice remain decisive. A survey can separate infrastructure that enables self-evolution from systems that actually perform it.

A fourth reason is safety. Self-evolution is attractive precisely because it can discover changes humans did not specify. That is also why it is risky. A self-modifying agent may remove a constraint, exploit a benchmark, learn to satisfy a proxy while violating the real objective, poison its own memory, or generate code with hidden vulnerabilities. A self-play system may collapse into trivial tasks. A reflection system may store incorrect lessons. A scientific automation system may flood review channels with plausible but shallow papers. These risks are not solved by refusing to build self-evolving systems; lightweight versions already exist in prompt optimizers, code agents, and memory systems. The community needs a shared framework for designing them safely.

A fifth reason is that self-evolution changes the unit of progress. The conventional unit is a model checkpoint. The self-evolutionary unit is a loop: an evolving population, an optimizer over prompts, a training curriculum generator, or an agent that patches itself. Progress may come from better verifiers, better archives, better rollback, better memory consolidation, better task proposal, or better governance, not only from larger base models. This shift matters for research strategy. If future capability gains come from systems that improve themselves around a foundation model, then benchmark reports must include the evolution protocol, compute budget, evaluator, and retention rule alongside the base model name.

\subsubsection{The role of open-source evidence and project health}

Open-source projects play two roles in this survey. First, they provide implementations through which self-evolution mechanisms can be inspected, reproduced, and extended. Second, they shape the language by which the broader community understands autonomy and agents. A framework with tens of thousands of stars can popularize a design pattern even if its internal self-evolution is limited. Conversely, a technically important research system may have modest adoption because it is difficult to run, expensive, or tied to a paper benchmark.

The star-analysis evidence used in this survey warns against equating popularity with quality. Hype-driven projects can accumulate stars rapidly through social media, demos, and narratives of autonomy. Academic or engineering-heavy projects often grow more slowly but attract higher-quality forks, issues, and contributors. The report's comparison among AutoGPT, Devika, CrewAI, MetaGPT, OpenEvolve, DSPy, SWE-Agent, OpenHands, LangGraph, AutoGen, FunSearch, and related projects motivates a more careful evaluation. We ask not ``Which project has the most stars?'' but ``Which project contains a feedback-driven modification loop, how reliable is that loop, and what evidence shows cumulative improvement?''

This distinction matters for the term ``Self Evolve'' as a brand and research area. A landing page, index, or community resource should not simply aggregate agent projects. It should identify where each project sits in the evolution stack. AutoGPT is historically important for popularizing autonomous task loops, but many versions relied on fixed prompting and tool use. CrewAI provides accessible role-based orchestration, but a fixed crew is not automatically self-evolving. MetaGPT encodes software-development roles, but its default pipeline is largely designed by humans. AutoGen supports agent conversation and human-in-the-loop control, which can host evolutionary loops but does not guarantee them. LangGraph offers a flexible state-graph substrate that can represent adaptive workflows. Letta/MemGPT focuses on memory evolution. DSPy offers prompt/program optimization. SWE-Agent and OpenHands operationalize software engineering with feedback from repositories and tests. OpenEvolve explicitly explores evolutionary code improvement. Each belongs in the ecosystem, but each should be labeled by mechanism.

Project health also affects scientific credibility. A self-evolving system is difficult to evaluate if the repository is abandoned, undocumented, or impossible to reproduce. Fork quality, issue quality, contributor diversity, and pull-request activity provide indirect evidence about whether a project can support cumulative research. The survey therefore treats open-source analysis as a complement to paper analysis. Papers provide claims and controlled experiments; repositories reveal engineering assumptions, tool dependencies, evaluation scripts, and community adoption. A mature self-evolution field will need both.

\subsubsection{Evaluation principles introduced in the survey}

Because self-evolving systems change over time, their evaluation must be temporal. We introduce several principles that recur throughout later chapters. The first is \emph{trajectory reporting}: papers should report not only the best final score but also the sequence of scores, failed modifications, accepted modifications, compute cost, and evaluation variance. An archive that finds one strong agent after hundreds of failed children is different from a stable optimizer that improves gradually. Both may be valuable, but they have different cost and risk profiles.

The second principle is \emph{held-out evolution}. A system can overfit not only a model but also an evolution process. If modifications are repeatedly selected on the same benchmark, the agent may learn benchmark-specific hacks. Therefore evaluation should distinguish development tasks used for evolution from held-out tasks used for final assessment. DGM's cross-domain transfer claims are important because they address whether evolved agents generalize beyond the search benchmark \cite{zhang2025darwin}. Similar held-out protocols are needed for prompt, memory, and curriculum evolution.

The third principle is \emph{verifier auditing}. In self-evolution, the verifier is a teacher, judge, and selection environment. If the verifier is wrong, the evolution process will optimize toward wrongness. Code executors are relatively strong but still incomplete if tests are weak. LLM judges are flexible but vulnerable to bias and inconsistency. Human feedback is valuable but expensive and variable. Reward models can be gamed. Verifier auditing should include adversarial examples, calibration checks, agreement with humans, and sensitivity to prompt changes.

The fourth principle is \emph{cost-normalized improvement}. A method that improves by one point after thousands of expensive model calls may be less useful than a method that improves slightly less with far lower cost. We therefore emphasize improvement per compute, per token, per environment step, or per wall-clock hour. AlphaEvolve-scale search may be justified for data-center scheduling or chip design because tiny percentage improvements have enormous value \cite{novikov2025alphaevolve}. The same cost would be unjustified for a low-stakes text-generation task.

The fifth principle is \emph{safety-preserving retention}. A candidate modification should not be retained solely because it improves a task score. It should also preserve constraints: security, privacy, tool permissions, alignment instructions, reproducibility, and rollback. This is especially important for code-level and architecture-level self-evolution. A patch that improves SWE-bench by disabling tests is not an improvement. A prompt that increases user satisfaction by ignoring safety policy is not an improvement. The acceptance function must encode both capability and constraints.

The sixth principle is \emph{diversity and novelty accounting}. Open-ended systems may discover useful stepping stones that do not immediately improve the main metric. Quality-diversity methods provide one way to preserve such variants. However, diversity should not become an excuse for retaining unsafe or irrelevant artifacts. The survey therefore treats diversity as a structured objective, ideally defined by behavioral descriptors or task niches rather than vague novelty.

\subsubsection{Safety and governance as part of the definition}

Safety is often treated as a later chapter in surveys, but for AI Self-Evolution it belongs in the definition. A self-evolving system without governance is not merely an incomplete product; it is an ill-specified optimizer. The acceptance function must decide what counts as an improvement. If it only measures benchmark score, the system may sacrifice interpretability, robustness, privacy, or compliance. If it only measures user approval, it may become sycophantic. If it only measures execution success, it may exploit tests. Therefore governance constraints $\mathcal{G}$ are part of the formal system tuple.

Governance begins with scope. A self-evolving agent should know which components it is allowed to modify. Prompt edits are lower risk than tool-permission changes; memory updates are lower risk than arbitrary code execution; local sandbox patches are lower risk than production deployment. The system should expose a clear modification surface. For example, an agent may be allowed to edit its task-specific prompt but not its safety preamble. It may be allowed to propose code patches but not apply them without tests and human approval. It may be allowed to generate training tasks but not change reward definitions. Scope limits turn self-evolution from an uncontrolled capability into an engineering pattern.

Governance also requires provenance. Every modification should be attributable: what feedback triggered it, what model proposed it, what files or components changed, what tests ran, what scores improved or regressed, and why it was accepted. Without provenance, debugging an evolving system becomes nearly impossible. Reflexion memories need timestamps and task links; prompt optimizers need version histories; code-evolving agents need diffs and test logs; archive-based systems need parent-child lineages. These records are not bureaucratic overhead. They are the empirical substrate for studying evolution.

Rollback is another governance primitive. A self-evolving system should be able to revert a harmful modification. G{\"o}del Agent explicitly includes rollback after failed validation \cite{yin2024godel}. Software-engineering practice offers many relevant tools: versioning, staged rollouts, canary deployments, continuous integration, feature flags, and audit trails. Future self-evolving agents should integrate these practices by default. The more powerful the modification surface, the stronger the rollback requirement.

Human oversight remains important, but it should be placed where it has leverage. Requiring humans to approve every minor reflection defeats the purpose of low-level self-improvement. Requiring no human approval for self-modifying production code is unsafe. A layered approach is preferable: automatic retention for low-risk memory entries, validation-gated retention for prompt and tool edits, sandbox-gated retention for code patches, and human approval for changes that affect permissions, safety policy, external side effects, or deployment. This layered governance mirrors the levels of self-evolution.

Finally, governance must account for multi-agent settings. If several agents can modify shared prompts, tools, memories, or code, conflicts and collusion become possible. One agent may optimize a component in a way that harms another. A critic may become too lenient toward a generator. A population may converge toward exploiting the same verifier. Multi-agent self-evolution therefore requires ownership, locking, review, and independent evaluation. These concerns are already familiar in human software teams; they become technical requirements when agents participate in their own improvement.

\subsubsection{A unifying vocabulary for later chapters}

Throughout the survey, we use several terms consistently. A \emph{substrate} is the representation being modified: text, memory, prompt, tool schema, graph, code, task distribution, weights, or population. A \emph{variation operator} proposes changes to that substrate. A \emph{critic} interprets performance and produces feedback. A \emph{verifier} provides a more objective signal, such as tests or environment rewards. A \emph{selector} decides what to retain. An \emph{archive} stores variants or histories. A \emph{governor} enforces constraints. An \emph{evolution trace} is the sequence of states, feedback, modifications, and accept/reject decisions over time.

These terms allow us to compare systems that otherwise appear unrelated. In Self-Refine, the substrate is an output; the variation operator is the same LLM; the critic is self-feedback; the selector may be a stopping rule. In Reflexion, the substrate is memory; the critic is a reflection model; the verifier is an evaluator; the selector appends feedback to memory. In Agent Symbolic Learning, the substrate is prompts/tools/pipelines; the variation operator is a language optimizer; the critic is a language loss; the selector updates symbolic weights. In ADAS, the substrate is agent code; the variation operator is a meta-agent; the verifier is a benchmark; the archive stores discovered agents. In DGM, the archive is central and parent selection becomes part of the evolutionary algorithm. In Absolute Zero, the substrate includes curriculum and weights; the verifier is a code executor; the selector is an RL update.

The vocabulary also helps diagnose missing components. A system with variation but no verifier will drift. A system with a verifier but no diversity mechanism may overfit. A system with an archive but no governance may retain unsafe variants. A system with strong critics but weak selectors may accumulate noisy memories. A system with strong selectors but weak variation may stagnate. Many current papers are best understood as improving one missing component: ReVeal improves verification; RAGEN improves trajectory-level RL and identifies echo traps; DGM improves archive-based code evolution; Agent Symbolic Learning improves symbolic credit assignment; Absolute Zero improves curriculum generation without external data.

\subsubsection{Limitations of this survey's scope}

Although the survey is broad, it does not claim that every adaptive AI system is self-evolving. We focus on systems where modification of operational components is explicit enough to analyze. This excludes much of ordinary supervised learning, one-off prompt engineering, static agent orchestration, and human-only model iteration. It also excludes speculative claims about recursive self-improvement that lack implemented feedback loops. Our aim is to build a grounded survey of existing mechanisms, not a catalog of science-fiction scenarios.

We also emphasize systems around LLMs and agents because the recent field is concentrated there. However, self-evolution is not intrinsically limited to language. Robotic policies can self-modify controllers, scientific agents can evolve experimental protocols, recommender systems can evolve objectives and exploration policies, and multimodal agents can adapt perception-action loops. The formal definition is representation-agnostic. The LLM focus reflects current evidence: language models are unusually capable variation operators over the artifacts that compose AI systems.

Another limitation is that many reported results are recent, preprint-based, or difficult to reproduce. Some systems are not fully open-source. Some benchmark numbers may change as evaluations improve. Some open-source project statistics are time-sensitive. We therefore use them as evidence for mechanisms rather than permanent rankings. The durable contribution is the framework: mutable substrate, feedback, modification, verification, retention, and governance.

Finally, we do not assume that self-evolution is always desirable. For many applications, static systems are safer, cheaper, and easier to certify. A medical triage model, financial decision system, or legal assistant may require strict version control and human approval. Self-evolution is most attractive where tasks are diverse, feedback is reliable, and improvement can be sandboxed before deployment. The survey's goal is not to promote uncontrolled self-modification, but to understand when and how feedback-driven system improvement can be made useful and safe.


\subsubsection{Design dimensions for comparing self-evolving systems}

The mechanisms surveyed in this paper can be compared along a set of design dimensions. The first dimension is \emph{self-access}. A system may have no access to its internals, read-only access, proposal access, sandboxed write access, or production write access. Prompt-level methods usually have write access only to a local context. Memory agents have write access to a memory store. Code-evolving systems may read and write source files in a sandbox. Runtime self-modifying agents may patch live methods. Weight-level systems update parameters through a training loop. The degree of self-access determines both capability and risk. A system that can only append reflections is limited but relatively safe; a system that can rewrite its own tool permissions is powerful but dangerous.

The second dimension is \emph{feedback objectivity}. Feedback ranges from subjective self-critique to externally verified execution. Self-Refine uses model-generated critique, which is flexible but not always reliable \cite{madaan2023selfrefine}. Reflexion often uses an evaluator or environment success signal \cite{shinn2023reflexion}. SelfEvolve, ReVeal, Absolute Zero, and AlphaEvolve benefit from executable verification \cite{jiang2023selfevolve,jin2025reveal,zhao2025absolute,novikov2025alphaevolve}. AI Scientist uses novelty checks, experimental metrics, and LLM-based peer review, mixing objective and subjective signals \cite{lu2024ai}. Feedback objectivity does not guarantee usefulness---a unit test can be incomplete---but it strongly affects whether improvements accumulate.

The third dimension is \emph{temporal horizon}. Some loops operate within a single inference episode. Others operate across repeated attempts on one task, across a benchmark suite, across a deployment lifetime, or across generations of agents. The longer the horizon, the more important memory consolidation, archive management, and drift control become. A one-step critique can be evaluated by final answer quality. A month-long evolving agent must be evaluated by stability, transfer, cost, and safety incidents. Long horizons also create delayed-credit problems: a modification that looks neutral today may enable future improvements.

The fourth dimension is \emph{granularity of change}. An output edit is fine-grained and reversible. A prompt edit may affect many future tasks. A tool-schema change can alter action affordances. A workflow-graph change can alter control flow. A code patch can introduce arbitrary new behavior. A weight update can change latent capabilities in ways that are hard to localize. A population update changes the distribution of available agents. Granularity determines how much validation is necessary. Fine-grained changes can be accepted with local checks; coarse-grained changes require broader regression suites.

The fifth dimension is \emph{retention topology}. Some systems overwrite a single state. Others append memory. Others maintain version histories, branches, or archives. Greedy overwrite is simple but vulnerable to regressions. Append-only memory preserves experience but can bloat or contradict itself. Versioned prompts allow rollback and ablation. Archive-based evolution, as in DGM and AlphaEvolve, preserves multiple lineages \cite{zhang2025darwin,novikov2025alphaevolve}. Population structures are especially important for open-endedness because they prevent premature convergence, but they also make deployment choices harder: which variant should be used for which user or task?

The sixth dimension is \emph{credit assignment}. When a system improves, which component caused the improvement? In a simple Self-Refine loop, the answer may be the latest critique. In a multi-tool agent, improvement may result from a prompt, a tool call, a planner, a verifier, or their interaction. Agent Symbolic Learning explicitly tries to assign language gradients to nodes \cite{zhou2024agent}. ReVeal uses turn-level adaptive credit assignment \cite{jin2025reveal}. RAGEN decomposes trajectories into state, thinking, actions, and rewards \cite{wang2025ragen}. Credit assignment is not only a training issue; it is also an interpretability and governance issue. If we cannot identify which change helped, we cannot confidently retain, transfer, or audit it.

The seventh dimension is \emph{search pressure}. Some systems use weak pressure, such as a single self-critique. Others use strong pressure, such as many generations of selection on a benchmark. Strong pressure can produce large gains, but it also increases overfitting to the evaluator. This is a familiar lesson from evolutionary computation and AutoML. When the search process is powerful, the benchmark becomes part of the training environment. Therefore benchmark secrecy, held-out tests, adversarial evaluation, and transfer checks become more important.

The eighth dimension is \emph{human role}. Humans may provide objectives, approve changes, label feedback, design verifiers, inspect logs, or intervene after failures. A system with human approval can still be self-evolving if it proposes and validates its own modifications. Conversely, a system without humans is not necessarily better; it may simply be less governed. We therefore treat human involvement as a design parameter rather than a disqualifying factor. The right level depends on risk and reversibility.

The ninth dimension is \emph{deployment coupling}. Some evolution happens offline in a research sandbox; some happens online during user interaction. Offline evolution allows controlled evaluation and rollback. Online evolution allows personalization and rapid adaptation but risks exposing users to unstable variants. Production self-evolution may require a two-stage structure: online logs generate candidate modifications, but acceptance occurs offline after validation. This structure mirrors continuous integration in software engineering and can prevent unsafe immediate self-rewrite.

The tenth dimension is \emph{model dependence}. Many self-evolution loops depend on frontier LLMs for variation quality. Agent Symbolic Learning notes the need for strong backbones for language-gradient computation \cite{zhou2024agent}. ADAS and DGM depend on foundation models to propose meaningful agent code \cite{hu2024automated,zhang2025darwin}. AlphaEvolve uses a dual-model strategy for breadth and depth \cite{novikov2025alphaevolve}. A loop that works only with the strongest proprietary models has different reproducibility and cost properties from a loop that works with smaller open models. Future research should report sensitivity to base-model strength.

\subsubsection{Illustrative formal decomposition of the evolution loop}

A compact way to compare systems is to decompose every self-evolution episode into six functions:
\begin{equation}
(S_t, x_t) \xrightarrow{\mathrm{act}} \tau_t \xrightarrow{\mathrm{judge}} F_t \xrightarrow{\mathrm{diagnose}} G_t \xrightarrow{\mathrm{propose}} \Delta_t \xrightarrow{\mathrm{validate}} V_t \xrightarrow{\mathrm{retain}} S_{t+1}.
\end{equation}
Here $\tau_t$ is the behavioral trace, $F_t$ is feedback, $G_t$ is a diagnosis or gradient-like explanation, $\Delta_t$ is a candidate modification, and $V_t$ is validation evidence. Different papers instantiate or omit different functions. Self-Refine merges judge and diagnose into self-feedback and may not use separate validation. Reflexion separates evaluator feedback from reflective diagnosis. Agent Symbolic Learning makes diagnosis explicit as language gradients. SelfEvolve uses interpreter errors as feedback and revised code as the candidate modification. DGM uses benchmark evaluation as validation and archive insertion as retention. ReVeal trains the validate step itself. Absolute Zero uses task proposal and solving to generate both $x_t$ and $\tau_t$.

This decomposition helps identify where progress is needed. If acting is weak, the system needs a better base model or tools. If judging is weak, it needs a better evaluator. If diagnosis is weak, feedback will not translate into useful changes. If proposal is weak, modifications will be syntactically or semantically poor. If validation is weak, harmful changes will be retained. If retention is weak, improvements will be lost or regressions will accumulate. A system is only as strong as the weakest function in this chain.

The decomposition also clarifies why language models are unusually useful. They can participate in every function. They can act by solving tasks, judge by critiquing outputs, diagnose by explaining failures, propose by editing prompts or code, validate by generating tests, and retain by summarizing lessons. But using the same model for all functions can create correlated errors. The model may fail to notice its own mistakes, produce a flattering diagnosis, propose a patch that satisfies its own flawed judge, and then retain the patch. Strong systems therefore diversify roles: external executors, separate evaluator models, human review, held-out tests, and archive-level selection all reduce correlated failure.

One can also view the evolution loop as approximate Bayesian or evolutionary inference. The system maintains a belief over useful modifications; feedback updates that belief; proposals sample from high-promise regions; validation estimates fitness; retention updates the population or state. This view suggests uncertainty-aware self-evolution. Instead of greedily accepting the top-scoring modification, a system might maintain confidence intervals, explore uncertain variants, or require stronger evidence for high-risk changes. Such methods are underdeveloped in current LLM-agent work, which often relies on small benchmark samples and deterministic accept/reject rules.

\subsubsection{From benchmark improvement to capability accumulation}

A central question is whether self-evolution produces durable capability accumulation or merely benchmark-specific improvement. A prompt optimized for one dataset may fail elsewhere. A code patch may pass public tests but fail hidden cases. A self-generated curriculum may teach narrow tricks. An archive may accumulate variants that are individually strong but hard to compose. Therefore the survey distinguishes \emph{score gain} from \emph{capability gain}. Score gain is measured on the evaluation used during evolution. Capability gain transfers to held-out tasks, new environments, or future evolution steps.

Capability accumulation requires reusable structure. Reflexion memories accumulate if they encode general lessons rather than task-specific reminders. Agent Symbolic Learning accumulates if prompt and tool updates improve underlying procedures. ADAS accumulates if discovered architectures transfer across domains or backbones. DGM accumulates if archive variants provide stepping stones for future agents. AlphaEvolve accumulates if evolved algorithms reveal patterns reusable across problem families. Absolute Zero accumulates if self-generated tasks induce reasoning skills that transfer beyond code. AI Scientist accumulates if generated research ideas and experimental strategies improve future scientific work rather than only producing isolated papers.

This distinction also affects how we interpret negative results. A self-evolving system may fail to improve a final score while still discovering useful components, diagnostics, or benchmarks. Conversely, it may improve a score while damaging generality. For that reason, later chapters advocate richer reporting: ablations of retained modifications, transfer tests, qualitative analysis of discovered strategies, and lineage diagrams. The goal is to know not only that $S_T$ is better than $S_0$, but why it is better and whether the reason will matter elsewhere.

The capability-accumulation lens connects self-evolution to open-endedness. In open-ended systems, the objective is not only to solve a known benchmark but to keep generating new possibilities. Quality-diversity archives, novelty metrics, environment generation, and self-play curricula are all mechanisms for preserving the capacity to discover. The challenge is to keep open-endedness productive. Unconstrained novelty can drift into irrelevant artifacts; pure objective pressure can converge prematurely. The balance between novelty and utility is one of the core design problems inherited from evolutionary computation.

\subsubsection{Implications for building self-evolving systems}

For practitioners, the framework suggests a staged approach. The safest starting point is usually a local reflection or verification loop around a fixed agent. For example, a code assistant can generate tests, run them, and revise code before responding. A research assistant can critique a summary against source documents. A customer-support agent can store approved corrections in memory. These loops provide immediate value while keeping the modification surface narrow.

The next stage is symbolic component optimization. Prompts, retrieval queries, tool descriptions, routing policies, and workflow graphs can be versioned and tuned against validation tasks. This stage should use explicit development sets, regression tests, and rollback. It is often more practical than weight fine-tuning because the components are interpretable and cheap to edit. Agent Symbolic Learning and DSPy-like systems point in this direction \cite{zhou2024agent}.

A further stage is sandboxed code evolution. Agents can propose patches to their own tools, planners, or evaluators, but patches should be isolated, tested, and reviewed. This stage benefits from software-engineering infrastructure: unit tests, static analysis, dependency scanning, permission boundaries, and CI. DGM and AlphaEvolve show the upside of code evolution, while G{\"o}del Agent shows the need for runtime safeguards \cite{zhang2025darwin,novikov2025alphaevolve,yin2024godel}.

The most advanced stage is population or weight-level evolution. Here systems generate tasks, update policies, maintain archives, or evolve multiple agents. This stage requires substantial compute and strong evaluation. It is suitable when the domain has reliable rewards or verifiers and when long-term improvement justifies the cost. Absolute Zero, RISE, RAGEN, ReVeal, DGM, and AlphaEvolve occupy different points in this stage \cite{zhao2025absolute,qu2024rise,wang2025ragen,jin2025reveal,zhang2025darwin,novikov2025alphaevolve}.

Across all stages, the practical advice is the same: make the loop explicit. Define the mutable substrate. Define the feedback. Define the proposal mechanism. Define validation. Define retention. Define rollback. Define logging. Define who can approve high-risk changes. A self-evolving system that cannot answer these questions is not an engineering system; it is an uncontrolled experiment.

\subsubsection{Research questions opened by the definition}

The definition introduced in this chapter leads to several research questions. First, what are the minimal conditions under which self-evolution improves generalization rather than overfitting? This question requires theory and benchmarks that model repeated selection on limited feedback. Second, how should language-gradient methods be formalized? Natural-language diagnoses are expressive but hard to aggregate, compose, or verify. Third, how can we design verifiers for domains where correctness is not executable? Many valuable tasks involve judgment, creativity, social context, or long-horizon consequences.

Fourth, how should archives be structured for agent evolution? DGM and AlphaEvolve show the value of retaining diverse variants, but archive size, sampling, pruning, niche definition, and safety filtering remain open. Fifth, how can self-evolving systems avoid reward hacking and evaluator exploitation? This problem is familiar in RL but becomes more acute when the system can modify code, prompts, or tests. Sixth, how can self-generated curricula remain challenging and aligned? Absolute Zero suggests one path, but task proposal is itself an optimization problem.

Seventh, how should we certify systems that change after deployment? Traditional certification assumes a fixed artifact. Self-evolution requires certifying a process, including its constraints and monitoring. Eighth, how should human oversight be allocated? Too much oversight prevents useful adaptation; too little invites unsafe drift. Ninth, what is the relationship between self-evolution and interpretability? Evolution traces may provide explanations of how a system changed, but weight-level updates may remain opaque. Tenth, how do multiple self-evolving agents interact when they share tools, memories, or environments?

These questions motivate the rest of the survey. They also show why AI Self-Evolution is not a narrow subtopic of prompting or AutoML. It is a systems problem that spans representations, algorithms, software infrastructure, evaluation science, and governance.


A final practical implication is that self-evolution should be reported with enough detail to be reproduced as a process. A paper should name the base model, prompts, tools, mutable files or components, feedback sources, evolution budget, acceptance thresholds, rollback policy, random seeds, benchmark splits, and all retained variants. Without these details, readers can observe only the final artifact and must guess the evolutionary path that produced it. This is analogous to reporting a trained model without the data or optimizer. Because the loop is the object of study, the loop must be documented. We therefore treat evolution traces, audit logs, and lineage records as first-class scientific artifacts throughout the survey.

This process view also changes how the field should discuss failure. Failed modifications are not merely noise to be hidden; they reveal the search landscape, evaluator weaknesses, and constraints that shape future designs. A system that fails safely and records why it failed may be more scientifically valuable than a system that reports only cherry-picked successes. For self-evolving AI to mature, negative evidence, regressions, and rejected patches should be shared alongside successful improvements.

\subsection{Survey Methodology}
\label{subsec:intro-methodology}

This survey follows a systematic literature collection and classification protocol adapted from PRISMA guidelines for scoping reviews \cite{moher2009preferred}. The methodology consists of four stages: source identification, screening, eligibility assessment, and classification.

\paragraph{Source identification.} We searched five electronic databases (arXiv, Semantic Scholar, Google Scholar, ACL Anthology, and NeurIPS/ICML/ICLR proceedings) using the query string: \texttt{("self-evolving" OR "self-improving" OR "self-modifying" OR "self-refining" OR "recursive self-improvement") AND ("AI agent" OR "LLM" OR "language model" OR "agentic system")}. The search covered publications from January 2022 to May 2026, capturing the emergence of the field after the widespread deployment of instruction-tuned LLMs. We additionally performed backward snowball sampling from included papers and forward citation tracking via Semantic Scholar. To complement academic sources, we collected 486 GitHub repositories tagged with agent, self-evolution, or self-improving keywords, yielding 204 analyzed projects with model-card documentation.

\paragraph{Screening.} After removing duplicates, we screened 327 candidate papers against the following inclusion criteria: (1) the system modifies one or more of its own operational components (prompts, memory, tools, code, architecture, data curriculum, or weights); (2) modification is driven by feedback from interaction, evaluation, or verification; (3) the system retains or accumulates modifications across episodes; and (4) sufficient technical detail is available to classify the evolution mechanism. We excluded papers that describe only static prompting, conventional fine-tuning with fixed datasets, or agent orchestration without feedback-driven self-modification.

\paragraph{Eligibility and classification.} Eligible papers were classified along five dimensions: mutable component (what evolves), feedback source (what drives change), verification mechanism (what validates change), retention topology (how change persists), and temporal scope (when evolution occurs). Each paper received at least two independent classifications by the authors, with disagreements resolved by discussion. The resulting taxonomy produces the Five Evolution Loops framework described in Section~\ref{subsec:intro-contributions}. We further classified open-source projects by their self-evolution coverage: projects that implement feedback-driven modification loops were labeled as ``core evolution,'' projects that host or enable such loops were labeled as ``evolution infrastructure,'' and projects that provide agent capabilities without feedback-driven modification were labeled as ``static orchestration.''

\paragraph{Evidence synthesis.} Quantitative evidence was extracted from benchmark results reported in original papers, with attention to evaluation protocol, sample size, base model, and reproducibility. Qualitative evidence was drawn from architecture analysis, code inspection, and community feedback (Hacker News, Reddit, X/Twitter discussions). We do not treat benchmark numbers as permanent rankings; instead, we use them as evidence for mechanism effectiveness under specific conditions. All classification data and project metadata are maintained in a versioned repository to support updates and corrections.

\subsection{Relation to Existing Surveys}
\label{subsec:intro-survey-comparison}

Several surveys touch on aspects of AI self-evolution, but none provides the cross-domain unification and mechanism-level taxonomy offered here. We distinguish our work along four dimensions.

\paragraph{Self-evolving LLM agents.} Wang et al.\ \cite{wang2025selfevolving} survey self-evolving LLM-based agents with a focus on feedback-driven improvement loops. Their taxonomy organizes methods by evolution stage (experience acquisition, refinement, updating), while our Five Evolution Loops framework classifies by the mechanism through which feedback becomes retained change. Our survey additionally covers evolutionary computation, AutoML/NAS, and production frameworks that their LLM-centric scope excludes. CharlesQ9/Self-Evolving-Agents \cite{charlesq9selfevolving} provides a GitHub-indexed resource list but without formal analysis or cross-domain comparison. YoungDubbyDu/LLM-Agent-Optimization \cite{youngdubbydu2026optimization} surveys optimization methods for LLM agents, focusing on prompt and pipeline tuning rather than the full spectrum of mutable components.

\paragraph{AutoML and neural architecture search.} Classical AutoML surveys \cite{he2021automl,yao2018taking} cover automated pipeline construction and hyperparameter optimization but do not address post-deployment self-modification, agent architectures, or code-level evolution. NAS surveys \cite{elsken2019neural,ren2021comprehensive} focus on neural network topology search rather than agent software, prompts, or multi-modal architectures. Our survey bridges AutoML/NAS with agent research by treating ADAS \cite{hu2024automated} and DGM \cite{zhang2025darwin} as descendants of NAS applied to agent programs.

\paragraph{Evolutionary computation and LLMs.} Surveys on LLM-augmented evolutionary algorithms \cite{yang2024large,liu2024large} cover LLMs as variation operators within evolutionary frameworks but do not address self-referential agent systems, reflection memory, or production deployment concerns. LLM4EC \cite{wuxingyu2025llm4ec} focuses on the intersection of LLMs and evolutionary computation for combinatorial optimization. Our survey treats evolutionary computation as one of several contributing traditions, alongside RL, program synthesis, and agent reasoning, unified under the self-evolution abstraction.

\paragraph{LLM agents and reasoning.} Recent surveys on LLM agents \cite{xi2025riseagent,wang2024surveyagent} cover tool use, planning, and multi-agent coordination but classify methods by capability rather than by evolution mechanism. Surveys on self-improving reasoning \cite{huang2025selfimprove} address chain-of-thought refinement but do not cover code-level or architecture-level self-modification. Our survey's distinguishing contribution is the mechanism-level taxonomy that compares methods across all seven mutable component types (prompts, memory, tools, code, architecture, data, weights) and the explicit treatment of safety, governance, and production deployment as integral to the self-evolution definition.

\begin{table}[t]
\centering
{\TableTight
\begin{tabularx}{\textwidth}{@{}L{0.22\textwidth}cccc@{}}
\toprule
\textbf{Dimension} & \textbf{This survey} & \textbf{Wang et al.\ 2025} & \textbf{AutoML surveys} & \textbf{EC+LLM surveys} \\
\midrule
Mutable components & 7 types & 3--4 types & 1--2 types & 1 type \\
Evolution mechanisms & 5 loops & Stage-based & Pipeline-based & Operator-based \\
Post-deployment self-modification & Yes & Partial & No & No \\
Agent architectures & Yes & Yes & No & No \\
Code-level evolution & Yes & Limited & No & Yes \\
Production frameworks & Yes & No & No & No \\
Safety and governance & Core definition & Separate section & Minimal & Minimal \\
Cross-domain unification & AutoML+NAS+EC+RL+agents & LLM agents only & AutoML only & EC only \\
\bottomrule
\end{tabularx}}
\caption{Comparison of this survey with related surveys along key dimensions. Our survey uniquely covers all seven mutable component types, five evolution loop mechanisms, production frameworks, and treats safety as part of the self-evolution definition.}
\label{tab:survey-comparison}
\end{table}

\subsection{Paper Structure}
\label{subsec:intro-structure}

The remainder of the survey is organized as follows. Chapter~2 develops a taxonomy of AI Self-Evolution. It expands the levels introduced here---output, memory, prompt/tool, architecture, code, curriculum, weights, and population---and provides a multidimensional classification by autonomy, feedback source, verification strength, retention mechanism, and safety boundary. The chapter also distinguishes closed-loop adaptation from open-ended evolution and introduces terminology for comparing systems that operate at different layers.

Chapter~3 studies core methods. It covers reflexive prompting, verbal reinforcement learning, textual backpropagation, prompt and tool optimization, multi-agent feedback, self-debugging, and verification-guided program synthesis. It treats Self-Refine, Reflexion, Agent Symbolic Learning, SelfEvolve, ReVeal, RISE, and RAGEN as method families rather than isolated papers. The chapter emphasizes how feedback is represented and transformed into modifications.

Chapter~4 focuses on evolutionary and open-ended systems. It connects NEAT, novelty search, quality diversity, POET, MAP-Elites, FunSearch, ADAS, DGM, AlphaEvolve, OpenEvolve, and AI Scientist. The chapter explains how LLMs change the role of mutation operators, how archives support cumulative discovery, and why open-endedness requires different evaluation principles from single-task optimization.

Chapter~5 addresses evaluation. It argues that self-evolving systems must be evaluated as trajectories. Standard pass@k, accuracy, win rate, or reward metrics capture endpoint performance but not evolvability, retention, cost, safety, transfer, or robustness. The chapter reviews benchmarks such as HumanEval, MBPP, DS-1000, SWE-bench, LiveCodeBench, HotPotQA, MATH, ALFWorld, WebShop, ARC, GFootball, and AI-scientist peer-review settings. It also discusses evaluator overfitting, benchmark leakage, reward hacking, echo traps, and reproducibility.

Chapter~6 surveys frameworks and implementation ecosystems. It compares agent frameworks and open-source projects such as AutoGen, LangGraph, Letta/MemGPT, DSPy, SWE-Agent, OpenHands, CrewAI, MetaGPT, AutoGPT, smolagents, OpenEvolve, ADAS, DGM implementations, Self-Refine repositories, Reflexion implementations, and AI Scientist. The chapter separates static orchestration from genuine self-evolution and analyzes project health using signals beyond star counts.

Chapter~7 studies pain points and risks. It covers unsafe self-modification, tool misuse, sandbox escape, evaluator gaming, reward hacking, memory poisoning, prompt drift, catastrophic forgetting, cost explosion, context-window limits, provenance, accountability, and governance. It also discusses why formal guarantees remain difficult and how practical systems can use rollback, audit logs, staged deployment, human approval, and constrained action spaces.

Chapter~8 outlines future directions. It proposes benchmark suites for self-evolution trajectories, standardized logging formats for evolution histories, safer self-modifying code protocols, hybrid formal-empirical verification, multi-agent evolutionary ecosystems, self-evolving scientific workflows, and governance standards for systems that can alter their own behavior. The chapter returns to the central thesis: future AI progress will depend not only on larger static models but also on reliable mechanisms by which AI systems improve themselves.

In summary, this introduction frames AI Self-Evolution as a practical and theoretical shift. The field inherits Turing's behavioral evaluation, Schmidhuber's self-referential ambition, AutoML's automated design, evolutionary computation's variation and selection, reinforcement learning's feedback optimization, and the LLM revolution's ability to manipulate language and code. Its defining object is the feedback-driven loop that modifies the system itself. The central questions are therefore no longer only ``How capable is the model?'' or ``How good is the agent?'' but also ``How does the system change itself, who verifies the change, what is retained, what is forgotten, what can go wrong, and how do we measure progress over an evolving trajectory?''
